In this section, we investigate examples and counterexamples for the various concepts defined in this article.
We hope this section highlights the non-trivial nature of finitely adequate modules.
So far, the most effective method for finding counterexamples is not to construct a specific module that is not finitely adequate, but instead to show that, within a given family of modules, not all the members are finitely adequate.

We begin by discussing various infinitesimally linear and tangential modules.
These properties are weaker than finite adequacy, and they serve as a good starting point to build intuition for the kinds of modules we are interested in.

\subsection{Infinitesimally Linear and Tangential Modules}
In this subsection, we discuss various infinitesimally linear and tangential modules, as well as modules that fail to satisfy these properties.
We also show that the implication from tangentiality to infinitesimal linearity is strict.
Schemes provide a natural starting point for this discussion.

\begin{example}\label{ex:sch-inf-lin}
  By \cite[Lemma 2.2.5]{diffgeo-article}, any scheme is infinitesimally linear at all of its points.
  In particular, the finite free $R$-modules $R^n$ are infinitesimally linear at any point.
\end{example}

\begin{example}\label{ex:R-lin-flat}
  The $R$-modules $R^n$ are tangential \cite[Lemma 2.2.3]{diffgeo-article}.
\end{example}

\begin{remark}
  If $M$ is both an $R$-module and a scheme, then $\tangent{M}{0}$ is linearly adequate by \cite[Lemma 2.3.8]{diffgeo-article}.
  In particular, if such an $M$ is tangential, then it is linearly adequate.
\end{remark}

Consider \Cref{prp:inf-lin-hood}, from which we can extract the following principle: if $M$ and $N$ are modules such that $\hood{M,0}{1} \subseteq N \subseteq M$, then $(M,0)$ is infinitesimally linear at $0$ if and only if $(N,0)$ is.

\begin{example}
  As an example, $\mathrm{Nil}(R)$ is infinitesimally linear at $0$ since it contains $\hood{R,0}{1} = \D(1)$.
\end{example}

On the other hand, the above principle behaves very differently when we consider tangential modules.
Suppose $M$ and $N$ are $R$-modules such that $\hood{M,0}{1} \subset N \subset M$.
Then, by \Cref{prp:hood-factorization}, we have $\tangent{M}{0} = \tangent{N}{0}$.
Thus, at most one of $M$ and $N$ can be tangential.

\begin{counterexample}\label{cex:nil-not-tangential}
  The nilradical $\Nil(R)$ is not tangential.
  There is a chain of proper inclusions $\D(1) \subset \Nil(R) \subset R$.
  Therefore, $\Nil(R)$ cannot be tangential, since $R$ is.
\end{counterexample}

\begin{counterexample}\label{cex:hood-p-not-tangential}
  Suppose that $R$ is a ring with positive characteristic $p > 0$. 
  Then we have a chain of proper inclusions $\hood{R,0}{1} \subset \hood{R,0}{p-1} \subset R$. 
  Therefore, $\hood{R,0}{p-1}$ cannot be tangential, since $R$ is.
\end{counterexample}

\begin{remark}
  By \Cref{cex:nil-not-tangential}, we see that being tangential is a strictly stronger condition than being infinitesimally linear.
\end{remark}

So far, all modules we have considered are infinitesimally linear.
However, not every module satisfies this condition; we present a counterexample below.

\begin{counterexample}
  The following quotient of free modules is not infinitesimally linear at $0$:
  \[
    M \defeq R[\D(2)]/R[\D(1) \vee \D(1)]
  \]
  To see why this is not infinitesimally linear at $0$, let $\eta : \D(2) \to R[\D(2)]$ be the unit map and $p : R[\D(2)] \to M$ the projection map.
  The composite map $p \circ \eta : \D(2) \to M$ is pointed and satisfies $(p \circ \eta)|_{\{0\} \vee \D(1)} = 0$ and $(p \circ \eta)|_{\D(1) \vee \{0\}} = 0$.
  However, since $\D(1) \vee \D(1) \to \D(2)$ is provably not surjective, we see that $p \circ \eta \neq 0$.
  Thus, the canonical map $(M,0)^{\D(2)} \to (M,0)^{\D(1)} \times (M,0)^{\D(1)}$ is not injective.
\end{counterexample}

So far, the examples of tangential modules we have constructed are all submodules of finitely presented modules.
We will now consider some bigger examples, in the sense that they are not sub-finitely presented.

\begin{example}\label{ex:power-series-tangential}
  By \Cref{prp:exp-tangential}, the modules $R[x] = R^R$ and $R[[x]] = R^\mathbb{N}$ are tangential.
  By \Cref{cor:tangential-inf-lin}, these modules are also infinitesimally linear.
  Note that $R[[x]]$ is not adequate, as it is not weakly quasi-coherent.
\end{example}

\subsection{Finitely Adequate Modules}

In this subsection, we present examples of finitely adequate modules, as well as counterexamples.
We are particularly interested in finding counterexamples among modules that satisfy several of the following conditions: being weakly quasi-coherent, infinitesimally linear, tangential, reflexive, or not-not finite free.
We begin by investigating finitely generated ideals of $R$.

\begin{example}\label{ex:fin-gen-ideal}
  Let $I$ be a finitely generated ideal of $R$.
  Then $I$ is finitely adequate.
  Indeed, if $r_1, \dots, r_n$ generate $I$, then $I$ is the image of the map $R^n \to R$ given by $(x_1, \dots, x_n) \mapsto \sum_{i} r_i x_i$. The result then follows from \Cref{thm:fin-adeq-class}.
\end{example}

We can construct interesting examples and counterexamples using three families of propositions.
For each $r \in R$, we define the propositions $V(r) \defeq (r = 0)$, $D(r) \defeq (r \neq 0)$, and $\neg D(r) \defeq \neg (r \neq 0)$.
Although one could also consider finite sequences of elements $r_1, \dots, r_n$ in $R$ and their joint propositions, it is sufficient to consider a single element to obtain our main counterexamples.

Let $U$ be a proposition.
The module $R^U$ is tangential by \Cref{prp:exp-tangential}, but it is not always finitely adequate.
We analyze $R^U$ as $U$ ranges over the three families defined above.
Observe that if $U$ is either $V(r)$ or $D(r)$, then $R^U$ is a finitely presented algebra by \cite[Definition 4.1.1]{draft} and \cite[Proposition 3.1.6]{draft}; consequently, these modules are quasi-coherent.
For the first case, $U = V(r)$, we obtain an example of a finitely adequate module.

\begin{example}\label{ex:exp-Vr-fin-adeq}
  Let $r : R$.
  The module $R^{V(r)} = \quot{R}{\ideal{r}}$ is finitely adequate by \Cref{ex:fin-gen-ideal} and \Cref{prp:fin-adeq-cok-closure}.
\end{example}

For the second family, $U = D(r)$, the module is the localization $R^{D(r)} = R[r^{-1}]$ by \cite[Proposition 2.2.3]{draft}.
Rather than constructing a specific counterexample, the next result shows that $R[r^{-1}]$ cannot be finitely adequate for all $r : R$.

\begin{counterexample}\label{cex:exp-Dr-not-fin-adeq}
  It is not the case that $R^{D(r)}$ is finitely adequate for all $r : R$.
  For the sake of contradiction, suppose that $R^{D(r)}$ is finitely adequate for all $r : R$.
  Then
  \[
    R[r^{-1}]^\ast = \{ a=(a_0,a_1,\dots) : R^{\N} \;|\; \forall k \geq 1.\; a_{k-1} = ra_k \}
  \]
  is also finitely adequate, and is thus weakly quasi-coherent.
  We now have a commutative square
  \[
  \begin{tikzcd}
    (R[r^{-1}]^\ast)[r^{-1}] \ar[d, "\simeq"]\ar[r] & R[r^{-1}], \ar[d, "\simeq"] &[-25pt] \frac{a}{r^n} \mapsto \frac{a_0}{r^n} \\
    (R[r^{-1}]^\ast)^{D(r)} \ar[r, "\simeq"] & R^{D(r)}, & f \mapsto f_0
  \end{tikzcd}
  \]
  where the lower horizontal map is an equivalence with inverse
  \[
    b \mapsto (p \mapsto (b(p),r^{-1}b(p),\dots))\rlap{.}
  \]
  This is well-defined because $p$ witnesses that $r$ is invertible.
  Accordingly, the upper horizontal map is an equivalence.
  The preimage of $\frac{1}{1}:R[r^{-1}]$ along this map is represented by an element $\frac{a}{r^n}$ such that $r^k(a_0-r^n)=r^k(r^{n+1}a_{n+1}-r^n)=r^{k+n}(ra_{n+1}-1)=0$ for some $k$.
  Since $R$ is local, we have $(ra_{n+1} \neq 0) \lor (ra_{n+1}-1 \neq 0)$, which readily implies that $(r \neq 0) \lor (r^{k+n} = 0)$.
  But then any element $r:R$ would either be invertible or nilpotent, contradicting the fact that $R$ is connected.
\end{counterexample}

For the third family, we consider $R^{\neg D(r)}$.
First, observe that, by \cite[Proposition 2.2.3]{draft}, $\neg D(r)$ is the proposition that $r$ is nilpotent.
Equivalently, $\neg D(r) = \bigcup_{n : \N} (r^n = 0)$. 
Using continuity of the exponential and duality, we obtain $R^{\lnot D(r)} = R^{\bigcup_{k : \N}(r^k = 0)} = \lim_k R^{(r^k = 0)} = \lim_k \quot{R}{\ideal{r^k}} = \widehat{R_r}$; i.e., $R^{\neg D(r)}$ is the $r$-adic completion of $R$.
If $r$ is either invertible or nilpotent, then $\widehat{R_r}$ is finitely adequate.
Our next counterexample follows the same strategy as the preceding one.

\begin{counterexample}\label{cex:exp-not-Dr-wqc}
  It is not the case that the module $R^{\lnot D(r)}$ is weakly quasi-coherent for every $r : R$.
  For the sake of contradiction, assume that $R^{\lnot D(r)}$ is weakly quasi-coherent for all $r : R$.
  By virtue of duality, the global sections of the bundle ${r \mapsto R^{\lnot D(r)}}$ form the power series module:
  \[
    R[[x]] = \lim_n R[x]/\ideal{x^n} = R^{\mathrm{colim}_n \Spec R[x]/\ideal{x^n}} = R^{\{r:R | \lnot\lnot(r=0)\}} = \prod_{r : R} R^{\lnot D(r)}\rlap{.}
  \]
  By \cite[Lemma 7.1.10]{draft}, we can deduce that $R[[x]]$ is weakly quasi-coherent since it is the global sections of a family of weakly quasi-coherent modules.
  However, $R[[x]]$ is not weakly quasi-coherent by \cite[Proposition A.0.5]{draft}, which is a contradiction.
  Thus, it follows that $R^{\lnot D(r)}$ is not weakly quasi-coherent for all $r$, and so it cannot be finitely adequate either by \Cref{prp:wqc-adeq}.
\end{counterexample}

For every proposition $U$, we define the type $R_U \defeq \sigmatype{x : R} (x = 0) \lor U$.
We first show that these modules are tangential, and for this we need a lemma.

\begin{lemma}\label{lem:forall-distributes-over-or}
Let $X$ be a projective type such that every map to $\Bool$ is constant.
Let $P, Q : X \to \Prop$ be predicates.
Then $\forall x : X, P(x) \lor Q(x)$ is equivalent to $(\forall x, P(x)) \lor (\forall x, Q(x))$.
\end{lemma}

\begin{proof}
Since the two types are propositions, it suffices to give implications in both directions.
There is an obvious implication from right to left, so we are left to prove that
$\forall x : X, P(x) \lor Q(x)$ implies $(\forall x, P(x)) \lor (\forall x, Q(x))$.
So assume $\forall x : X, P(x) \lor Q(x)$.
Since $\lor$ is the propositional truncation of the coproduct, the fact that $X$ is projective
means that we merely have a map in $\prod_x P(x) + Q(x)$.
Since we are proving a proposition, we can drop the ``merely''.
This gives us a map $X \to \Bool$, which by assumption must be constant.
This implies that we have either $\prod_x P(x)$ or $\prod_x Q(x)$, as required.
\end{proof}

\begin{proposition}\label{prp:R_U-tangential}
For any proposition $U$, the modules $R_U$ and $R/R_U$ are tangential.
\end{proposition}

\begin{proof}
We first consider $R_U \defeq \sum_{x : R} (x = 0) \lor U$.
The type $(R_U,0)^{\D(n)}$ is equivalent to
\[
  \sum_{f : \D(n) \rightarrow_\ast R} \forall \epsilon : \D(n), (f(\epsilon) = 0) \lor U .
\]
Since $\D(n)$ is projective (\Cref{lem:disk-projective}) and all maps $\D(n) \to \Bool$
are constant (\cite[Proposition 1.3.4]{diffgeo-article}), we can apply \Cref{lem:forall-distributes-over-or} to
get an equivalence to
\[
  \sum_{f : \D(n) \rightarrow_\ast R} (f = 0) \lor U ,
\]
where we have replaced $\forall \epsilon : \D(n), f(\epsilon) = 0$ with $f = 0$
and have used that $\D(n)$ is pointed to replace $\D(n) \to U$ with $U$.
Since $R^n$ is tangential, we can replace $\D(n) \rightarrow_\ast R$ with $R^n$,
giving
\[
  \sum_{x : R^n} (x = 0) \lor U ,
\]
which equals $(R_U)^n$.
It is not hard to check that the resulting equivalence is homotopic to the canonical map $(R_U)^n \to (R_U,0)^{\D(n)}$.
For $n=1$, this implies that $R_U$ is tangential, and considering the commutativity of the diagram
\[
  \begin{tikzcd}
    (R_U)^{n+m} \ar[r, "\simeq"]\ar[d, "id"] & (R_U,0)^{\D(n+m)} \ar[d] \\
    (R_U)^n \times (R_U)^m \ar[r, "\simeq"] & (R_U,0)^{\D(n)} \times (R_U,0)^{\D(m)}
  \end{tikzcd}
\]
we see that $R_U$ is infinitesimally linear and hence tangential.
It follows from \Cref{prp:op-closure-inf-lin-tangential} that $R/R_U$ is tangential as well.
\end{proof}

% This is almost entirely due to Claude, with a bit of guidance from jdc.
% The writing has been edited.
\begin{counterexample}\label{cex:R_Dr-not-wqc}
  We want to show that $R_{D(r)}$ can not be weakly quasi-coherent for all $r : R$.
  Indeed, if it were, then the quotient $\quot{R}{R_{D(r)}}$ would also be weakly quasi-coherent, since $R$ is weakly quasi-coherent and the quotient of a weakly quasi-coherent module by a weakly quasi-coherent submodule is weakly quasi-coherent \cite[Proposition 7.1.13]{draft}.
  Thus, it is sufficient to show that it can not be the case that $\quot{R}{R_{D(r)}}$ is weakly quasi-coherent for all $r : R$.

  Now, for the sake of contradiction, assume that $R/R_{D(r)}$ is weakly quasi-coherent for every $r : R$.
  Fix $r : R$, and recall that $D(r) \defeq (r \neq 0)$ is a proposition and that
  $R_{D(r)} = \sum_{x : R} (x = 0) \lor (r \neq 0)$.

  By weak quasi-coherence, the canonical map
  $(R/R_{D(r)})[1/r] \to (R/R_{D(r)})^{D(r)}$ is an equivalence by \Cref{def:quasi-coherent}.
  We claim that $(R/R_{D(r)})^{D(r)} = 0$.  Indeed, if $r \neq 0$, then the
  predicate $(x = 0) \lor (r \neq 0)$ holds for every $x : R$, so $R_{D(r)} = R$ and
  hence $R/R_{D(r)} = 0$.  This implies $(r \neq 0) \to (R/R_{D(r)} = 0)$,
  and so every element of $(R/R_{D(r)})^{D(r)}$ is zero.

  It follows that $(R/R_{D(r)})[1/r] = 0$.
  In particular, the class of $1$ is annihilated by a power of
  $r$: there merely exists $k : \N$ with $r^k \in R_{D(r)}$, that is,
  $(r^k = 0) \lor (r \neq 0)$.
  This is equivalent to $(\exists k.\ r^k = 0) \lor (r \neq 0)$.
  Since any nilpotent element satisfies $\lnot\lnot(r = 0)$, and $r \neq 0$ is equivalent to $\lnot(r = 0)$, we obtain $\lnot\lnot(r = 0) \lor \lnot(r = 0)$.
  As $r : R$ was arbitrary, we deduce that
  $\prod_{r : R} \big(\lnot\lnot(r = 0) \lor \lnot(r = 0)\big)$, a form of weak excluded
  middle.  The two disjuncts are mutually exclusive,
  so the truncation in $\lor$ can be dropped.  Recording which disjunct holds defines
  a function $\chi : R \to \Bool$, which takes different values at $r = 0$ and $r = 1$.
  Thus $\chi$ is not constant,
  which contradicts the connectedness of $R$ \cite[Proposition 1.2.2]{topology-draft}.
\end{counterexample}

\begin{remark}\label{rmk:R^U-R_U-not-not-fin-adeq}
  Let $U$ be a proposition.
  The modules $R^U$ and $R_U$, which provide the counterexamples considered above, are not-not finitely adequate.
  Indeed, these modules are finitely adequate when $U$ is either true or false.
  Since $\lnot\lnot(U \lor \lnot U)$ holds, the claim follows by double-negation elimination.
\end{remark}

We conclude this section by presenting examples of reflexive modules that are not finitely adequate.

\begin{counterexample}\label{cex:polynomial-power-series-not-fin-adeq}
  Neither $R[x]$ nor $R[[x]]$ is finitely adequate.
  In fact, they are not finite free, and thus cannot be not-not finite free, which is required for finitely adequate modules.
\end{counterexample}

\begin{lemma}\label{lem:perfect-pairing-powerseries-polynomials}
  The pairing
  \[
    R[x] \times R[[x]] \to R, \;\;
    (a,b) \mapsto \langle a,b \rangle :\equiv \sum_j a_jb_j
  \]
  is perfect.
\end{lemma}

For the proof, we present a constructive and adapted version of the one given in \cite[Satz 1]{Specker1950}.
This is a folklore theorem due to Ernst Specker, and the proof in loc.\ cit.\ shows, in particular, that $\Z[x]$ is isomorphic to $\Hom_{\Z}(\Z^\N, \Z)$.
While Specker's proof uses prime factorization in $\Z$, we replace this with a neat degree-counting argument using duality.

\begin{proof}
  It is clear that the map $\vartheta' : R[[x]] \to \Hom_R(R[x],R)$ defined by $b \mapsto (a \mapsto \langle a,b \rangle)$ is an equivalence. Thus, we only need to show that the second induced map
  \[
    \vartheta:
    R[x] \to \Hom_R(R[[x]],R), \;\;
    a \mapsto \left(b \mapsto \langle a,b \rangle\right)
  \]
  is an equivalence.
  Since $\vartheta$ is clearly injective, we only need to verify its surjectivity.
  Let $\lambda : R[[x]] \to R$ be an arbitrary $R$-linear map and define a sequence $a \defeq (n \mapsto \lambda(e^n))$, where $e^n$ is the sequence given by $j \mapsto \delta_{nj}$.
  We will show that $a_n = 0$ for $n \gg 0$, and then finish the proof by verifying $\lambda = \langle a, \argmhole \rangle$.
  Define further a map
  \[
    f : R \to R, \;\;
    r \mapsto \lambda(1,r,r^2,r^3,\dots)\rlap{,}
  \]
  Duality tells us that $f$ is a polynomial.
  Since we are proving a proposition, we may assume that $\deg f \leq d$ for some $d$.
  By the linearity of $\lambda$, for any $N : \N$, we have
  \[
    f(r) = a_0 + a_1r + a_2r^2 + \dots + a_{N-1}r^{N-1} + r^N\lambda(0,\dots,0,1,r,r^2,\dots)\rlap{,}
  \]
  where exactly the first $N$ entries in $\lambda(0,\dots,0,1,r,r^2,\dots)$ are zero.
  Comparing coefficients, we see that $a_N=0$ for $N \geq d+2$, hence $a$ is an element of $R[x]$.
  So it remains to check that $\varphi : \equiv \lambda - \langle a,- \rangle = 0$.
  To this end, let $b : R[[x]]$ be arbitrary and set
  \[
    g:R \to R, \;\;
    r \mapsto \varphi(b_0,b_1r,b_2r^2,\dots)\rlap{.}
  \]
  Since $\varphi(e^n) = 0$, we can repeat the argument just applied to $f$ and see that all polynomials $x^n$ divide $g$.
  But then $g=0$, and in particular $0 = g(1) = \varphi(b) = \lambda(b) - \langle a,b \rangle$.
\end{proof}

\begin{corollary}
  The $R$-modules $R[[x]]$ and $R[x]$ are reflexive.
\end{corollary}

Since $R[[x]]$ is not finitely adequate, we see that not every reflexive $R$-module is finitely adequate.

\begin{proof}
  We have a commutative diagram
  \[
  \begin{tikzcd}
    R[[x]]\arrow{r}{\ev_{R[[x]]}} \arrow{dr}{\vartheta'} & R[[x]]^{\ast\ast} \arrow{d}{\vartheta^\ast}
    \\
    & R[x]^\ast
  \end{tikzcd}
  \]
  where $\vartheta$ and $\vartheta'$ are as in the proof of \Cref{lem:perfect-pairing-powerseries-polynomials}.
  As shown there, these two maps are equivalences, which implies that $\ev_{R[[x]]}$ is as well.
  Consequently, $R[x]$ is reflexive by \Cref{cor:dbl-dual-reflexive-iss-dual-reflexive}.
\end{proof}


\subsection{Tensor products}

In this subsection, we show that the tensor product of two linearly adequate modules can fail to be linearly adequate.
This shows that the double dual tensor is the appropriate left adjoint to consider for the internal hom of $\finadeq$ (\Cref{cor:dbl-dual-tensor-adjunction}).

\begin{lemma}
  \label{tensor-gamma-commute-on-affines}
  Let $X$ be an affine scheme and suppose that $\mathcal F, \mathcal G : X \to \Mod R$ are weakly quasi-coherent bundles.
  Then the canonical map
  \[
    \Gamma(X,\mathcal F) \otimes_{R^X}\Gamma(X,\mathcal G) \to \Gamma(X,\mathcal F \otimes \mathcal G), \;\;
    f \otimes g \mapsto (x \mapsto f_x \otimes g_x)
  \]
  is an $R^X$-linear isomorphism.
\end{lemma}

\begin{proof}
  \rednote{TODO.}
\end{proof}

\begin{counterexample}\label{cex:tensor-not-closed}
  Finitely adequate modules are not closed under tensor products.
  Indeed, let $\varepsilon : \D(1)$ be arbitrary and consider the annihilator $M_\varepsilon :\equiv \ker(R\overset{\cdot \varepsilon}{\to}R)$.
  We show that $N_\varepsilon :\equiv M_\varepsilon \otimes M_\varepsilon$ cannot be reflexive for all $\varepsilon : \D(1)$.

  Since dualizing is exact, we have $M_\varepsilon^\ast = R/(\varepsilon)$.
  Then we have
  \[
    N_\varepsilon^\ast = \Hom_R(M_\varepsilon,R/(\varepsilon)) = \coker(M_\varepsilon^\ast \overset{\cdot \varepsilon}{\to} M_\varepsilon^\ast) = R/(\varepsilon)\rlap{,}
  \]
  where, for the penultimate equivalence, we use that $\Hom_R(M_\varepsilon,\argmhole)$ is right exact by \Cref{identification-hom-tensor-for-fp-modules}.
  In particular, $N_\varepsilon^{\ast\ast} = M_\varepsilon$, and under this identification the evaluation map $\ev : N_\varepsilon \to M_\varepsilon$ is the multiplication $s \otimes t \mapsto st$.
  Hence $\ev(\varepsilon\otimes\varepsilon)=0$ and thus even $\varepsilon \otimes \varepsilon = 0$ by virtue of our assumption.
  Now, using \Cref{tensor-gamma-commute-on-affines}, compute
  \[
    \Gamma(\D(1),N) = \Gamma(\D(1),M) \otimes_A \Gamma(\D(1),M) = I \otimes_A I = I \otimes_R I \rlap{,}
  \]
  where $A :\equiv R[x]/\ideal{x^2}$ and $I :\equiv \ker(A \overset{\cdot x}{\to} A) = \ideal{x}/\ideal{x^2}$.
  Note that the zero-section $(\varepsilon \mapsto \varepsilon \otimes \varepsilon) : \Gamma(\D(1),N)$ corresponds to $x \otimes x$ under these maps.
  As this element is not zero, we obtain a contradiction.
\end{counterexample}
